\section{Formal Verification of Reference Orbits}
\label{sec:reference-orbit-proofs}

One high-precision reference supplies the orbit samples used across a
deep-zoom view. Perturbation evolves the offsets from that reference for
individual pixels, so a misplaced reference sample can affect many pixel
calculations (\cref{sec:perturbation-concept,subsec:ref-orbit-pipeline}).
The Lean proofs connect the CPU implementations to the quadratic recurrence,
then follow the effects of rounding through the samples each implementation
publishes. For the fixed ST, MT3, and MT9 backends, the resulting bounds
relate the MPIR limb count to the spacing of the calculation grid.
For dynamic MT9, a bound on the additional drift shows when a reduced
working precision still meets the same viewport tolerance.

\paragraph{Notation.}
The reference parameter \(c_\star=c_{\star,x}+i\,c_{\star,y}\), exact
reference orbit \(z_n^\star=x_n^\star+i\,y_n^\star\), derivative \(d_n\),
and reduced-precision shadow \(\hat z_n\) follow
\cref{tab:notation-general,tab:notation-components,tab:notation-reference-orbit}.
The index \(n\) counts mathematical iterates from \(z_0^\star=0\), and \(j\)
counts loop bodies from zero. The fixed-precision proofs use the following
notation for the internal orbit and its errors.

\begin{center}
\begin{tabularx}{\linewidth}{@{}l>{\raggedright\arraybackslash}X@{}}
\(z_n^{\mathrm{MPIR}},\ \widetilde c_\star\) & Internal MPIR iterate and working
parameter. At a selected sample, \(z_n^{\mathrm{MPIR}}\) supplies the
high-precision value \(z_n^{\mathrm{HP}}\) before conversion to \(\hat z_n\). \\
\(\lVert z\rVert_1,\ s_n\) & Componentwise norm
\(\lVert z\rVert_1=|\Re z|+|\Im z|\) and state size
\(s_n=\lVert z_n^{\mathrm{MPIR}}\rVert_1\). \\
\(P,\ u\) & Effective MPIR precision in bits and error unit \(u=2^{-P}\).
The symbol \(p\) denotes the period. \\
\(\varepsilon_n,\ \delta\) & State error
\(\varepsilon_n=z_n^{\mathrm{MPIR}}-z_n^\star\) and parameter discrepancy
\(\delta=\lVert\widetilde c_\star-c_\star\rVert_1\). \\
\(\rho_n,\ \tau_n\) & Local arithmetic residual
\(\rho_n=z_{n+1}^{\mathrm{MPIR}}-((z_n^{\mathrm{MPIR}})^2+\widetilde c_\star)\)
and its bound in the componentwise norm. \\
\(E_n^{\mathrm{orb}}\) & Bound on \(\lVert\varepsilon_n\rVert_1\).
For MT9, \(F_n\) denotes the calculator's bound for fixed precision and
\(E_n\) denotes the Euclidean radius enclosing the difference between
dynamic and fixed precision. \\
\end{tabularx}
\end{center}

The symbols \(\varepsilon_n\), \(\rho_n\), and \(\tau_n\) serve the same roles
in the MT9 analysis. The replay error \(e_n\) concerns reconstruction from
stored orbit samples, and the pixel offset \(\Delta c=c-c_\star\) concerns
the position of a pixel relative to the reference parameter.

The standalone library is in \code{FractalSharkLean}. With Elan installed,
run \code{lake build} from that directory. The toolchain file and Lake manifest pin
Lean and Mathlib. A first build needs network access to resolve dependencies.
\code{lake exe cache get} downloads precompiled Mathlib artifacts. Lake keeps
build products in the package's \code{.lake} directory. The default build
checks the examples and reports the axiom dependencies of the main theorems.
The umbrella import is \code{FractalSharkProofs}. The limb model is also
available through \code{FractalSharkProofs.MPIR}.

\subsection{MPIR arithmetic assumptions}
\label{subsec:reference-orbit-proof-mpir}

All three CPU backends rely on the same limb representation. The primitive
contracts establish which operations preserve a value exactly and how much
rounding an operation can introduce. They supply the starting bounds for
every CPU precision result in this section.

The Lean model describes the
\href{https://github.com/BrianGladman/mpir/tree/dc82b0475dea84d5338356e49176c40be03a5bdf}{MPIR revision pinned by the Windows x64 build}.
For signed mantissa $m$, used limb count $\ell$, and MPIR exponent $e$, the
represented value is $mB^{e-\ell}$ with limb radix $B=2^{64}$. If
$t=\code{\_mp\_prec}$, then the object has room for $t+1$ limbs, reports
$P=64(t-1)$ bits, and yields the error unit $u=2^{-P}$.
The correspondence to C is established by source inspection. It assumes valid
initialized objects, supported aliasing, exact low-level limb arithmetic, and
exponent and size calculations within range. The working parameter, live
coordinates, and scratch objects use a common default precision.

\begin{itemize}
\item Initialization yields zero. Assignment and non-aliased negation retain
      up to $t+1$ limbs. In-place negation changes the sign. Swap exchanges
      the values and their precision metadata.
\item Whole-limb shifts retain up to $t+1$ input limbs. Other power-of-two
      shifts retain $t$ input limbs and reserve space for a carry.
      Multiplication retains up to $t$ limbs from each input and $t+1$
      limbs of the normalized product.
\item Same-sign addition aligns both operands at a shared exponent, discards
      low limbs, and retains the carry. Subtraction accounts for common
      leading limbs, borrow chains through zero and maximal limbs, and
      normalization of a surviving tail. The theorem quantifies over the
      limb boundary selected by the source branch.
\end{itemize}

The assignment error is bounded by $u|a|$, and equal-precision assignment
preserves the value. The sum and difference errors are bounded by
$2u\max(|a|,|b|)$, the product error by $3u|ab|$, and the error of a
power-of-two shift by $u|2^k a|$. The Lean routine contracts are reusable
independently of the orbit proofs.
The theorem \code{MPIR.add\_aligned\_bounds} also gives exact addition
when both represented values are on the retained limb grid, including
operands stored with trailing zero limbs.

\subsection{Single-threaded exact orbit}
\label{subsec:reference-orbit-proof-st-exact}

The serial implementation provides a baseline for the optimized backends.
Identifying its samples by mathematical iteration gives the perturbation
loop the correct reference term even though the first loop body starts
with the parameter already loaded.

The serial update in \cref{subsec:ref-orbit-st} evaluates the two components as
\[
 (x_n^\star+y_n^\star)(x_n^\star-y_n^\star)+c_{\star,x},
 \qquad (2x_n^\star)y_n^\star+c_{\star,y}.
\]
The theorem \code{SingleThreaded.exactStep\_eq} establishes that both
expressions read the same old state and give
\eqref{eq:mandelbrot-recurrence} in exact arithmetic.
\code{SingleThreaded.exactPrefix\_eq} identifies the synthetic zero entry
with $z_0^\star$, the initial loop state with $z_1^\star=c_\star$, and
sample $j+1$ with the state observed before the update in loop body $j$.

\subsection{How rounding errors accumulate}
\label{subsec:reference-orbit-proof-propagation}

Deeper views require finer absolute accuracy. A small error introduced by
a parameter copy or an early arithmetic operation can grow over many
subsequent iterations, especially when the orbit spends time at states with
large magnitude. The CPU precision proofs track that accumulated error at
each mathematical index \(n\).
Subtracting the exact recurrence from the computed update gives
\begin{equation}
  \varepsilon_{n+1}
  = (z_n^{\mathrm{MPIR}}+z_n^\star)\varepsilon_n
      +(\widetilde c_\star-c_\star)+\rho_n.
  \label{eq:reference-orbit-proof-error-identity}
\end{equation}
The product is the effect of squaring a state that already contains error.
Its factor follows from the difference of the squares of
\(z_n^{\mathrm{MPIR}}\) and \(z_n^\star\).
The parameter discrepancy enters every step because both maps add their own
parameter. The local residual accounts for the new rounding in the update
from \(z_n^{\mathrm{MPIR}}\).

If \(\lVert\varepsilon_n\rVert_1\le E_n^{\mathrm{orb}}\), then \(z_n^\star\)
lies within \(E_n^{\mathrm{orb}}\) of \(z_n^{\mathrm{MPIR}}\) and
\(\lVert z_n^{\mathrm{MPIR}}+z_n^\star\rVert_1\le 2s_n+E_n^{\mathrm{orb}}\).
The componentwise norm bounds the magnitude of a product by the product of
the two norms. Together with \eqref{eq:reference-orbit-proof-error-identity},
this gives the recurrence checked by \code{prefixBound\_correct}.
\begin{equation}
\begin{aligned}
  E_0^{\mathrm{orb}} &= 0, \qquad E_1^{\mathrm{orb}} = \delta, \\
  E_{n+1}^{\mathrm{orb}} &=
    (2s_n+E_n^{\mathrm{orb}})E_n^{\mathrm{orb}}+\delta+\tau_n,
    \qquad n\ge1.
\end{aligned}
\label{eq:reference-orbit-proof-orbit-bound}
\end{equation}
The initial zero is exact and the first live state is \(\widetilde c_\star\).
The theorem \code{parameterCopy\_error} bounds the error from copying an
already represented parameter by \(u\) times its norm. If the input parameter
approximates \(c_\star\), its existing discrepancy also contributes to \(\delta\).

The bound controls the sum of the absolute coordinate errors of the internal
state. At fixed state magnitudes, each additional bit of \(P\) halves the
local rounding allowance. The parameter discrepancy \(\delta\) has its
own contribution. When \(2s_n+E_n^{\mathrm{orb}}\) exceeds one, existing error
can be amplified. The estimate allows all contributions to reinforce one
another, which makes it conservative.

At matched precision, parameter discrepancy, working-parameter norm, and
state magnitudes, MT3 has the smaller local allowance.
\begin{equation}
\begin{aligned}
  \tau_n^{\mathrm{ST}}
    &=u\bigl(32s_n^2+2\lVert\widetilde c_\star\rVert_1\bigr),\\
  \tau_n^{\mathrm{MT3}}
    &=u\bigl(9s_n^2+2\lVert\widetilde c_\star\rVert_1\bigr).
\end{aligned}
\label{eq:reference-orbit-proof-step-budgets}
\end{equation}
For fixed \(s_n\), the shared recurrence increases with its current
nonnegative bound and local allowance. It therefore gives
\(E_n^{\mathrm{orb}}\) for MT3 no larger than for ST under those matched
conditions. The difference comes from the \(s_n^2\) term. Each computed
run has its own rounded states and actual residuals.

\subsection{Single-threaded finite-precision orbit}
\label{subsec:reference-orbit-proof-st-precision}

The single-threaded bound makes the rounding cost of the factored formula
explicit. It lets precision selection use ST's actual operation order when
the renderer generates a serial reference.

The serial real update rounds the sum and difference before multiplying them.
The error of each factor therefore propagates into the product. The imaginary
update rounds the scaled operand before evaluating its product. The scratch
swap preserves the resulting value. Composing the primitive bounds for both
coordinates gives the ST allowance in \eqref{eq:reference-orbit-proof-step-budgets}.
\code{SingleThreaded.prefix\_error\_le} applies the shared propagation rule
to every index of the modeled prefix, and
\code{SingleThreaded.error\_at\_observed} specializes it to the observed samples.

\subsection{MT3 exact orbit and worker schedule}
\label{subsec:reference-orbit-proof-mt3-exact}

MT3 overlaps the two squares and reuses the imaginary square from its
lookahead. The proof connects that overlap to the iteration sequence
consumed by perturbation. In particular, it associates each published
sample with a coherent pair of coordinates.

\code{MT3.exactStep\_eq} proves that the two square workers and the caller
implement \eqref{eq:mandelbrot-recurrence} in exact arithmetic. The workers
form $(x_n^\star)^2$ and $(y_n^\star)^2$ from stable snapshots, and the caller forms the
cross product and assembles the components. The worker model associates
each square with its input snapshot and iteration index. Either worker may
finish first. The caller saves the square of the current imaginary component
before giving the Y worker the next imaginary component. The next dispatch
retains that lookahead job's index and operand. An already claimed job may
finish through a pending stop request before the caller joins the worker.

MT3 has no delayed sample to publish in body zero. Body $j\ge1$
publishes the sample associated with index $j$, while the workers operate on the
current snapshots. The serial schedule instead observes the sample at index $j+1$
in body $j$. The observation theorems cover the synthetic zero entry and
both sample sequences. On MT3 early exit, the terminal calculation may
write scratch derived from an earlier Y square. The proved prefix ends
at the last coherent sampled state.

\paragraph{Mailbox synchronization.}
The square, shadow, and input buffers are ordinary memory shared between
the caller and workers. The memory-order proof identifies the
synchronization that makes a worker's data visible before it is read or
reused on another thread.
The mailbox model relates the successful claims to the publications they
read. In the original ordering, the first Y input store and the worker's
result publication were relaxed, while the caller consumed the result with
release ordering. The model exhibits conflicting payload accesses without a
happens-before relation on those paths. In the authoritative X and Y square
workers, an input publication is release and the worker claims it with
acquire. The worker publishes the completed square and shadow with release,
and the caller consumes them with acquire. The paired operations order the
ordinary payload accesses. They also order the caller's reads and next
input writes before the worker reuses the buffer for the next job.
The source correspondence assumes that each successful claim reads the
matching publication. The logical model treats scheduling progress as a
condition for completion.

\subsection{MT3 finite-precision orbit}
\label{subsec:reference-orbit-proof-mt3-precision}

MT3's expanded update reaches the same exact recurrence as ST through a
different sequence of rounded operations. Its own local allowance lets the
same propagated-error calculation assess both backends in the same units.

The MT3 caller multiplies the two coordinates first and then doubles the
rounded product with \code{mpf\_mul\_ui}. In the Windows MPIR multiply-by-two
branch, the top bit of the highest discarded limb can contribute a carry
to the retained product. The primitive theorem accounts for that carry
and bounds the doubling error by $u|2a|$ for input $a$. The square
snapshots and the saved Y-square use equal-precision assignment. The
independently rounded squares feed a subtraction, so their rounding errors are
bounded directly by the two square magnitudes. Combining those bounds with the
cross-product and addition errors gives the MT3 allowance in
\eqref{eq:reference-orbit-proof-step-budgets}. The Lean theorem
\code{MT3.error\_at\_observed} uses that allowance for each coherent
transition in \eqref{eq:reference-orbit-proof-orbit-bound}.

\subsection{MT9 split products and worker handoff}
\label{subsec:reference-orbit-proof-mt9-exact}

The split products allow large multiplications to use several workers.
Exact reconstruction lets the rounding analysis charge only the final
stored product, and the buffer and ordering proofs justify assembling it
from the current batch's partial products before those buffers are reused.

The integer input of a split product is the full occupied MPIR mantissa.
The theorem \code{MT9.assembled\_product} proves the split and reconstruction
identity in \eqref{eq:mt9-karatsuba}. The theorem \code{MT9.batch\_fits}
proves that the low, high, and sum products fit the active buffer widths,
the shifted middle term fits the destination range, and the full product
leaves two zero guard limbs. The intermediate subtraction is nonnegative.
Padding and odd split widths preserve the input integers. These theorems
assume exact integer operations on their valid operand ranges.

The theorem \code{MT9.follower\_run\_before\_wait} gives the ordering from
the caller's release of each follower's job to the lead's wait for it.
With atomic write/read coherence, the lead may consume the current
completion once it observes an idle slot. The theorem
\code{MT9.current\_idle} excludes earlier completions under the
one-outstanding-batch invariant. Release/acquire publication of partial
products orders assembly, and the lead's completion orders the output
store. The next batch reuses snapshots and active widths only after
these stores. The same ordering protects aliased stores and shutdown
after a completed batch.

\subsection{MT9 exact and fixed-precision orbits}
\label{subsec:reference-orbit-proof-mt9-precision}

The fixed MT9 recurrence establishes the meaning and accuracy of the
samples produced by the split-product backend. It also supplies the
comparison trajectory for the dynamic backend at the same starting
precision.

The theorem \code{MT9.exact\_at\_observed} proves the exact recurrence at
all published indices. MT9 peels the initial check and update before
it publishes a positive index. Each later body checks the current
state and publishes the preceding state. A terminal body completes
its update before acting on periodicity or bailout. The theorem
\code{MT9.checked\_before\_publication} associates each positive
published index with its earlier check.

For finite arithmetic, \code{MT9.normalize\_product} accounts for
sign, actual occupied widths, and the exponent adjustment.
\code{MT9.product\_bounds} trims only the assembled product to the
output capacity. The same expanded update bound serves MT3 and MT9
with their respective product allowances. It gives
\[
 \lVert\rho_n^{\mathrm{F}}\rVert_1
 \le u\bigl(7s_n^2+2\lVert\widetilde c_\star\rVert_1\bigr).
\]
The theorem \code{MT9.error\_at\_observed} propagates that allowance
against the exact orbit. \code{MT9.viewportPrefix} supplies the
reference and ideal perturbation theorems with the fixed-precision
coefficient \(K_{\mathrm{MT9}}=7\).

\subsection{MT9 dynamic precision and the error disk}
\label{subsec:reference-orbit-proof-mt9-dynamic}

Using fewer limbs reduces the work in MT9's multiplications.
The question for rendering is how much that saving can change the
reference orbit. The comparison uses the orbit that fixed MT9 would
compute from the same copied center at the original precision.
MT9 Dynamic tracks a bound on the difference during its own run,
so a second orbit calculation is unnecessary.

At iteration \(n\), the theorem \code{MT9.disk\_prefix} proves that the
fixed state \(z_n^{\mathrm F}\) lies within distance \(E_n\) of the
dynamic state \(z_n^{\mathrm D}\) in the complex plane. The inequality
\(\lvert z_n^{\mathrm D}-z_n^{\mathrm F}\rvert\le E_n\) describes that disk.

The radius starts at zero because the calculations start from identical
values and use the same full-precision arithmetic until the first
reduction. From then on, each update accounts for any low limbs discarded
when precision drops, the effect of the existing difference on the next
squaring, and the new rounding in both calculations.
Returning to full precision carries forward the uncertainty already
accumulated. The proof applies to every permitted sequence of working
precisions under the limb and outward-rounding assumptions in
\cref{subsec:ref-orbit-dynamic-enclosure}.

The radius gives a bound on the extra error due to dynamic precision.
If the calculator bounds the fixed orbit's error against the exact orbit
by \(F_n\), then \(F_n+2E_n\) bounds the dynamic orbit's error in the
componentwise norm. The real and imaginary differences between the two
MT9 states are each at most \(E_n\). When the sum fits the viewport budget
at every iteration included in the reference output, MT9 Dynamic meets the
same tolerance at the original starting precision. A sum above the
budget leaves accuracy undecided. The worked example in
\cref{subsec:reference-orbit-proof-mt9-viewport} uses this test to show
that a run starting at 448 bits can drop to 384 bits and still meet
the quarter-sample tolerance.

\subsection{MT9 accepted conversion certificates}
\label{subsec:reference-orbit-proof-mt9-certificate}

Perturbation and compression consume converted samples. Equal conversions
therefore preserve their reference input even when the internal MPIR states
differ.

\code{MT9.regular\_gaps} subtracts the gap measurement error from
the conservative lower bounds used in the regular check.
\code{MT9.fast\_gaps} uses the leading 32 discarded bits and the
unread tail interval to lower-bound both gaps. Strict separation
keeps the fixed and dynamic components in one normalized bin with
the same sign. The theorem \code{MT9.sample\_equivalence} then gives
equality of the converted pair under the supported MPIR conversion
contract. A zero disk gives equality directly, and a real parameter
gives an independent exact-zero theorem for the imaginary component.
An unsupported format, an invalid bound, or an ambiguous gap leaves
the check unaccepted. Failure records the first diagnostic while
the publication path continues.

\code{MT9.accepted\_prefix\_equivalence} uses the earlier check for
each delayed sample and includes the shared synthetic zero.
\code{MT9.same\_exit\_of\_sample\_equivalence} proves that agreeing
converted samples produce the same exit and period metadata in
the modeled loop when the center, viewport radius, shadow arithmetic,
and cancellation decisions agree. The proof covers the samples
entering storage or compression.

\subsection{A single tolerance check for a finite prefix}
\label{subsec:reference-orbit-proof-tolerance}

A single upper bound on the state sizes lets a precision decision cover an
entire finite reference prefix. The accumulated bound explains why a
precision that suffices for a short reference can become inconclusive for
the same view at a higher iteration count.

Let \(n_{\max}\) be the last observed orbit index to certify and \(T\) the
required absolute tolerance in the componentwise norm. Given an upper bound
\(S\) on \(s_n\) for \(1\le n<n_{\max}\), define a fixed amplification
\(A=2S+T\) and a per-step budget
\[
  q=\delta+u\bigl(KS^2+2\lVert\widetilde c_\star\rVert_1\bigr),
  \qquad K_{\mathrm{ST}}=32,\quad K_{\mathrm{MT3}}=9,\quad K_{\mathrm{MT9}}=7 .
\]
Here \(S\) bounds the internal orbit states. The view radius \(R_{\max}\)
used by the periodicity check describes the view geometry.

As long as the current error radius is at most \(T\), its amplification is
bounded by \(A\), and the parameter and local arithmetic terms sum to at
most \(q\). The Lean corollary \code{prefix\_error\_le\_of\_tolerance}
uses that fact to compare the true error to an affine bound \(F_n^{\mathrm{aff}}\).
\begin{equation}
\begin{aligned}
  F_0^{\mathrm{aff}}&=0, \qquad F_1^{\mathrm{aff}}=\delta, \qquad
  F_{n+1}^{\mathrm{aff}}=A F_n^{\mathrm{aff}}+q
       &&\text{for } n\ge1,\\
  F_{n_{\max}}^{\mathrm{aff}}
     &= A^{n_{\max}-1}\delta+
        q\sum_{k=0}^{n_{\max}-2} A^k
       &&\text{for } n_{\max}\ge1 .
\end{aligned}
\label{eq:reference-orbit-proof-tolerance}
\end{equation}
The sum has \(n_{\max}-1\) terms and is empty at \(n_{\max}=1\).
For \(A>1\), it grows geometrically with the prefix length. At \(A=1\),
it grows linearly as \(n_{\max}-1\). For \(0\le A<1\), it is bounded by
\(1/(1-A)\). For any \(A\ne1\), it equals \((A^{n_{\max}-1}-1)/(A-1)\).

Since \(A\ge0\) and \(q\ge\delta\ge0\), \(F_n^{\mathrm{aff}}\) is nondecreasing.
A single check \(F_{n_{\max}}^{\mathrm{aff}}\le T\) therefore bounds every earlier radius by \(T\)
and justifies the error induction. The two
\code{error\_at\_observed\_of\_tolerance} theorems apply this result to ST and
MT3 with their respective coefficients and prove
\(\lVert\varepsilon_n\rVert_1\le T\) for every \(n\le n_{\max}\).
A result above tolerance leaves the requested accuracy undecided. The size
cap may follow from a separate orbit bound or from the computed prefix.
Numerical application requires exact
arithmetic or outward bounds for the inputs and the terminal calculation.

\paragraph{Worked comparison.}
Suppose \(S=1\) bounds both prefixes, with \(P=64\),
\(\lVert\widetilde c_\star\rVert_1=1\), \(\delta=0\), \(n_{\max}=40\), and
\(T=2^{-20}\approx9.54\times10^{-7}\). Then \(A=2+T\), and
\(q_{\mathrm{ST}}=34u\) while \(q_{\mathrm{MT3}}=11u\).
The MT3 bound is approximately \(3.28\times10^{-7}\), which meets
the tolerance for the whole prefix. The ST bound is approximately
\(1.01\times10^{-6}\), so the ST tolerance check is inconclusive.

\subsection{Viewport scale and CPU limb count}
\label{subsec:reference-orbit-proof-viewport}

Expressing error in sample spacings makes the tolerance depend on what
the renderer is asked to resolve. A deeper zoom, larger output, or denser
antialiasing grid can require more limbs for the same orbit.

The spacings \(h_x,h_y\) of the calculation grid come from the full
horizontal and vertical spans divided by \(sW\) and \(sH\), respectively,
where \(W\times H\) is the output size and \(s\) is the supersampling factor
per axis. See \cref{sec:mapping,subsec:ref-orbit-view-bounds} and
\eqref{eq:periodicity-pixel-spacing}. For a chosen fraction \(0<\eta\le1\),
set \(T=\eta\min(h_x,h_y)\). The theorem \code{Viewport.sampleError\_le}
shows that an orbit error \(\varepsilon\) with \(\lVert\varepsilon\rVert_1\le T\)
satisfies
\begin{equation}
  \frac{|\Re\varepsilon|}{h_x}+\frac{|\Im\varepsilon|}{h_y}\le\eta .
  \label{eq:reference-orbit-proof-sample-error}
\end{equation}
This comparison uses the parameter grid as a unit of length for the orbit
error. Image displacement and escape decisions involve additional dynamics.

\paragraph{An error bound for a given precision.}
Choose a fixed error ceiling \(\overline T\) and let \(C\) cap the states
whose errors are propagated. For the reference orbit, \(C=S\). Use the
amplification \(A=2C+\overline T\), and define \(G_m(A)=\sum_{k=0}^{m-1}A^k\)
for \(m=n_{\max}-1\) when \(n_{\max}\ge1\). Let \(c_{\mathrm{in}}\) be the
represented reference coordinate supplied to MPIR, \(M\) an upper bound
on \(\lVert c_{\mathrm{in}}\rVert_1\), and \(\beta\) an upper bound on
\(\lVert c_{\mathrm{in}}-c_\star\rVert_1\). The working parameter
\(\widetilde c_\star\) has norm at most \(M\). The coordinate copy contributes
at most \(uM\), so the discrepancy in \eqref{eq:reference-orbit-proof-orbit-bound}
is bounded by \(\beta+uM\).

Substituting these bounds into the affine formula gives the split checked
by \code{mpirViewportBound\_split}.
\begin{align}
  V_A &= A^m+G_m(A), \notag\\
  E_{\mathrm{input}} &= V_A\beta,
  & J &= V_A M+G_m(A)(KS^2+2M), \notag\\
  E &= E_{\mathrm{input}}+2^{-P}J.
  \label{eq:reference-orbit-proof-limb-budget}
\end{align}
At \(n_{\max}=0\) the bound is zero, and at \(n_{\max}=1\) only the parameter
contributes. When \(E\le\overline T\), the amplification used to obtain
\(E\) is valid throughout the prefix. The additional comparison \(E\le T\)
certifies \eqref{eq:reference-orbit-proof-sample-error} at each index.

With \(L=\min(\overline T,T)-E_{\mathrm{input}}\), the equivalent
precision check is \(J\le 2^P L\). For \(J>0\) and \(L>0\), the smallest
sufficient internal limb count is the first \(t\ge2\) for which
\(2^{64(t-1)}L\ge J\). If \(L\le0\), a positive \(J\) prevents this check
from passing at any finite precision. An extra limb divides the MPIR
contribution by \(2^{64}\) at fixed \(J\). The bound on the original input
must be established separately.

For a centered square-pixel view with vertical half-height \(R_{\max}\),
the grid spacing is \(2R_{\max}/(sH)\). The lower radius for the sample
budget is therefore \(sH E/(2\eta)\) at fixed caps and limb count.
The proof uses the half-height of the exact viewport, whose geometric role
corresponds to \code{GetMaxRadius()}.

\paragraph{How the current CPU request spends its guard bits.}
The heuristic in \cref{par:precision-heuristics} extracts exponents
\(e_x,e_y\) from the full spans via \code{HDRFloat<double>} and
\code{mpf\_get\_2exp\_d}. It requests \(b=d+g\), where
\(d=\max(|e_x|,|e_y|)\) and \(g\) is 120 bits in the standard path or
800 bits when authoritative reuse is requested. The viewport setter
installs the default used by \code{mpf\_init} in ST and MT3. In the
64-bit MPIR model, \(t=1+\lceil\max(53,b)/64\rceil\) is the internal
\code{\_mp\_prec}, with \(P=64(t-1)\) reported bits. For an overridden
default, the certificate uses the precision of the initialized buffers.

When the extracted exponents give lower span bounds \(2^{e_x-1}\) and
\(2^{e_y-1}\), both spans are at least \(2^{-d-1}\). Hence the smaller
pitch is at least \(2^{-d-1}/(s\max(W,H))\).
For \(\beta=0\), the theorem \code{Viewport.cpu\_roundoff\_le\_budget}
proves that the following sufficient guard-bit condition gives \(E\le T\).
\begin{equation}
  2s\max(W,H)J \le \eta\,2^{g+r},\qquad r=P-b .
  \label{eq:reference-orbit-proof-cpu-guard}
\end{equation}
The check \(E\le\overline T\) still applies. The guard budget covers grid
density, the requested fraction of a sample, and orbit amplification.
The precision-selection routine supplies the request, while the Lean
criterion uses the extracted exponent bounds and the modeled orbit
prefix to determine whether that request is sufficient.

\subsection{Viewport-wide ideal perturbation}
\label{subsec:reference-orbit-proof-ideal-perturbation}

The reference check concerns a single trajectory. The ideal perturbation
proof explains how arithmetic error in that reference can enter the
reconstructed trajectories of pixels throughout the view when their
offsets and updates are exact.

For each pixel parameter \(c\), evaluate the update from
\cref{sec:perturb-only-update} in exact arithmetic using
\(z_n^{\mathrm{MPIR}}\) as the reference and \(c-c_{\mathrm{in}}\) as
the offset. Denote that perturbation by \(\Delta z_n^{\mathrm{ideal}}\)
and the reconstructed state by
\(z_n^{\mathrm{rec}}=z_n^{\mathrm{MPIR}}+\Delta z_n^{\mathrm{ideal}}\).
The algebra checked by \code{idealReconstruction\_residual} yields
\begin{equation}
  z_{n+1}^{\mathrm{rec}}=(z_n^{\mathrm{rec}})^2+c+
    (\widetilde c_\star-c_{\mathrm{in}})+\rho_n .
  \label{eq:reference-orbit-proof-reconstructed-residual}
\end{equation}
The working-parameter copy error enters each step together with the
original MPIR residual. The comparison orbit uses the pixel parameter
\(c\) itself, so \(\beta=0\) in its budget.

The bounds \(R_x,R_y\) of \cref{subsec:ref-orbit-view-bounds}, taken relative
to \(c_{\mathrm{in}}\), give
\(\lVert c-c_{\mathrm{in}}\rVert_1\le R_x+R_y\) over the rectangle.
Given a perturbation cap \(D\ge0\), apply the geometric sum to the
perturbation recurrence. The check
\begin{equation}
  (R_x+R_y)\,G_{n_{\max}-1}(2S+D)\le D
  \label{eq:reference-orbit-proof-perturbation-cap}
\end{equation}
bounds every perturbation used as an input to a transition through
\(n_{\max}\ge1\). The zero prefix is exact. With \(C=S+D\) in
\eqref{eq:reference-orbit-proof-limb-budget}, the rounding term still
uses \(KS^2+2M\), since the MPIR operations are on the reference.
The Lean theorem \code{perturbed\_error\_in\_samples} therefore gives
\eqref{eq:reference-orbit-proof-sample-error} for every sample in the
rectangle when both the geometry and error checks pass.

For a centered square-pixel view about \(c_{\mathrm{in}}\),
\(R_x+R_y=(1+W/H)R_{\max}\). If \(n_{\max}\ge2\), the upper radius is
\(D/[(1+W/H)G_{n_{\max}-1}(2S+D)]\). The lower radius comes from the
sample budget. The resulting interval can be empty for the chosen
caps, precision, and prefix. The theorems use exact affine mapping and
unre-based perturbation on the observed internal states. Bounding the
rendered result also requires bounds for reduced samples, rebasing,
acceleration, and escape decisions.

\paragraph{Worked CPU comparison.}
Use the \(800\times400\) dimensions and \(s=2\) of \cref{sec:mapping}.
Take \(\eta=1/4\) and \(R_{\max}=2^{-73}\), so \(T=R_{\max}/1600\).
Suppose \(S=M=1\), \(\beta=0\), and the observed prefix reaches index 40.
Use \(D=\overline T=2^{-8}\). Three internal limbs (\(P=128\)) meet the
reference budget for both backends and the ideal perturbation budget
for MT3. Four limbs (\(P=192\)) also meet the ideal ST budget.

The full spans \(4R_{\max},2R_{\max}\) have extracted exponents \(-70,-71\).
The standard CPU request is 191 bits and rounds to four limbs. A reuse
request for the same spans is 871 bits and rounds to 896 bits, or 15
internal limbs. For a longer prefix through index 112, the same standard
request leaves the reference criterion inconclusive even with the same
size cap. The precision request depends on the spans, and the sufficient
bound also depends on the number of transitions.

\subsection{Selecting sufficient MPIR precision from a viewport}
\label{subsec:reference-orbit-proof-precision-estimator}

The calculator turns the error analysis into a precision request that
can be evaluated for a particular view. Its exact-orbit enclosures include
the amplification along that reference and give a sufficient working
precision for the specified prefix.

The standalone calculator selects the MPIR request before running the CPU
orbit. Run the examples from \code{FractalSharkLean}.
\begin{lstlisting}
lake exe orbit-precision examples/viewport-st.json
lake exe orbit-precision examples/viewport-st-zoom.json
lake exe orbit-precision examples/viewport-st-delayed-bailout.json
lake exe orbit-precision examples/viewport-st-periodicity.json
\end{lstlisting}
Specify \(c_\star\) with \code{centerX} and \code{centerY}. The viewport
accepts either \code{minX}, \code{maxX}, \code{minY}, and \code{maxY}, or
\code{zoomFactor}. These values accept exact decimals, including
scientific notation. In zoom mode, the three inputs correspond to Real,
Imaginary, and Zoom in Enter Location.

Enter Location starts with a square of half-height \(2/Z\) for zoom
\(Z>0\). Fitting the output aspect ratio gives the centered half-extents
\begin{equation}
  R_x=\frac{2W}{Z\min(W,H)},\qquad R_y=\frac{2H}{Z\min(W,H)} .
\end{equation}
The calculator uses exact rational arithmetic for the fit. The GUI
rounds its fitted bounds and can shift the reference center.

The calculation grid uses \code{width} and \code{height} for the output
size and \code{antialiasing} for the number of samples per axis. The
backend is \code{ST}, \code{MT3}, or fixed \code{MT9}, and \code{lastIndex} counts from
\(z_0^\star=0\). Optional \code{sampleFraction}, \code{bailoutRadius},
and \code{maxBits} default to \(1/4\), \(16\), and \(65{,}536\),
respectively.

The output gives separate settings for the reference and for ideal
perturbation, and a combined setting. Set MPIR's default precision to
\code{requestedBits} before initializing the orbit temporaries. Each
center component is encoded as a signed integer \code{mantissa} times
\(2^{\text{\code{binaryExponent}}}\). The mantissa fits in the destination
slot. It can be loaded by \code{mpf\_set\_z} and scaled by a
whole-limb power-of-two shift. Both assignments from the loaded center
to the working parameter and the initial state preserve the value.
The implementation checks that each encoded coordinate has magnitude
at most that of the requested coordinate and differs by at most
\(u\) times the requested magnitude, with \(u=2^{-P}\).

\paragraph{Deriving the precision request.}
The calculator encloses each exact reference state in a rectangle
using rational arithmetic and rounds the two component intervals
outward to a dyadic grid after every update. It tightens the grid
within the configured limit when interval widths exceed one.
Let \(U_n\) be the resulting bound on \(\lVert z_n^\star\rVert_1\),
\(M=\lVert c_\star\rVert_1\), and \(T=\eta\min(h_x,h_y)\) as in
\cref{subsec:reference-orbit-proof-viewport}. As long as the reference
error is at most \(T\), the computed state has norm at most \(U_n+T\).
The calculator propagates dimensionless error weights \(J_n\).
\begin{align}
  J_0 &= 0,\qquad J_1=M,\notag\\
  J_{n+1} &= \left\lceil (2U_n+T)J_n+
    M+K(U_n+T)^2+2M\right\rceil,\qquad n\ge1 .
  \label{eq:reference-orbit-precision-estimator}
\end{align}
The integer ceiling bounds the new weight while keeping the arithmetic
size manageable. If \(u\max_{0\le n\le n_{\max}}J_n\le T\), induction
proves the reference tolerance at every observed index. The factors
come from the supplied coordinate's orbit, and periodicity can shorten
the prefix covered by the same request.

For ideal perturbation, a second set of weights uses a reference
error ceiling of one to bound the reference arithmetic. The pixel
state cap at each step follows from its displacement from the exact
reference and the configured bailout radius. Ideal perturbation
uses exact offsets from the returned dyadic center, which is also
the working parameter. The parameter-copy term then cancels in
\eqref{eq:reference-orbit-proof-reconstructed-residual}. The resulting
guarantee covers the exact pixel prefix through its first bailout
iterate and accounts for the reference's MPIR error under ideal
perturbation. Quantization of stored samples, acceleration, and
escape-count equality require separate bounds.

The calculator selects the first supported bit count \(P\) satisfying
each bound by comparing the weights with \(2^P T\). The ideal setting
also has to fit its reference error ceiling, and the combined setting
satisfies both tests. The returned settings have the smallest limb
counts for these inequalities. The Lean theorems
\code{Recommendation.reference\_correct} and
\code{Recommendation.ideal\_correct} give the guarantees under the
MPIR and CPU prefix contracts already used in this section. A precision
or enclosure limit can leave a request without a certificate.

\subsubsection{Reference orbits near the left tip}
\label{subsec:reference-orbit-left-tip}

At \(c_\star=-2\), the exact reference orbit is
\(0,-2,2,2,\ldots\). All of its states and arithmetic intermediates are
small integers. The orbit lands on a repelling fixed point at \(2\),
where the derivative \(2z\) of \(z^2+c_\star\) equals \(4\).
A small error in the orbit state can therefore grow by about a
factor of four per step while the state remains close to \(2\).
That amplification costs roughly two bits of accuracy per step.

For a nearby real center \(c_\star=-2+\Delta c\), take
\(w_n=z_n^\star-2\), so \(\Delta c\) here measures the displacement of
the reference center from \(-2\). The first positive state is
\(z_2^\star=2-3\Delta c+(\Delta c)^2\). The recurrence becomes
\begin{equation}
  w_{n+1}=4w_n+w_n^2+\Delta c.
  \label{eq:reference-orbit-left-tip}
\end{equation}
A small \(|\Delta c|\) keeps the exact orbit near \(2\) for many
iterations. Early arithmetic errors can grow throughout the stay near
\(2\), and the fine sample spacing of a deep view can make low-order
terms such as \(w_n^2\) significant while the leading digits of
\(z_n^\star\) still barely change. The ST examples below exercise the
first effect. The reduced-precision MT9 example in
\cref{subsec:reference-orbit-proof-mt9-counterexample} shows how
discarded bits can exceed the tolerance on a single update.

\paragraph{Why the span-based margin can fall short.}
The examples distinguish an inconclusive upper bound from an actual
failure of the span-based precision request. They illustrate how
reference error can exceed the sample spacing and change the observed
reference bailout even when the rendered image looks normal.

The first two files specify a \(5120\times2160\) view with \(s=2\),
zoom \(2^{74}\), and \(c_\star=(-2+2^{-73})+0i\). The explicit
box rounds the horizontal bounds outward. Both inputs use an ST
prefix through index 80 with periodicity disabled. The standard
120-bit margin requests 191 bits from the spans. MPIR uses 192 effective
bits and four internal limbs. With the input coordinate loaded exactly,
the Windows x64 MPIR operations in ST order give an absolute error of
about \(6.606\times10^{-25}\) at index 80, as checked against an outward
enclosure of the exact orbit. The quarter-sample tolerance is
\(1.225\times10^{-26}\), so the error is about 54 times the tolerance.

The span-based request depends on the viewport size. The required
precision also depends on the time the orbit spends in the
amplifying region described by \eqref{eq:reference-orbit-left-tip}. The
calculator certifies the ST reference prefix at 256 bits, or five
internal limbs. The combined setting is 448 bits with the default
bailout radius.

\paragraph{A delayed reference bailout.}
The third file keeps the same viewport scale and sets
\(c_\star=(-2-5\cdot 2^{-193})+0i\). At index 99, the exact orbit
is about \(174.8643\). Four-limb Windows x64 MPIR with the entered
center loaded exactly gives about \(54.6165\), an error of \(120.25\).
The calculator recommends 320 bits for the reference prefix and
576 bits for the combined guarantee.

To reproduce either case in the GUI, use the HDRx64 GPU renderer
with LA~v2, a \(5120\times2160\) render area, and 4x antialiasing
(\(s=2\)). Select \emph{Single Thread (ST)} under \emph{Perturbation},
then use \emph{Clear Perturbation References - All}. Copy the center
and zoom from the JSON file into Enter Location and use an iteration
limit of 8192. Enter Location preserves the selected reference mode.
Check for \code{PerturbationAlg = ST} and \code{ActualPeriodIfAny = 0}
in the render details.

In the delayed-bailout case, the fitted center
is \(-2-3\cdot 2^{-193}\) on Windows x64. At index 99, its exact orbit
is about \(54.6165\), while the four-limb orbit reaches \(10.1677\).
The resulting error is 44.45, about \(3.63\times10^{27}\) times the
quarter-sample tolerance. The exact orbit crosses the ST reference
bailout threshold \(\lVert z_n+c\rVert_2>16\) at index 99. The computed
orbit crosses it at index 100. These measurements concern the internal
reference orbit. The rendered image can still look normal.

Auto selects \code{STPeriodicity} at the 191-bit coordinate request.
For this view, \(R_y=2^{-73}\) and \(z_n\) stays close to \(2\) in the
early iterations. The derivative grows until \(2R_y\lVert d_n\rVert_\infty\)
rises from about \(2/3\) at index 37 to \(8/3\) at index 38, triggering
the near-return test of \cref{subsec:ref-orbit-periodicity}. The orbit
contains only the 38 computed entries and the initial \(z_0=0\) entry.
The stop is a heuristic for a nearby critical return, and the exact
orbit at the fitted center escapes.

\paragraph{A precision failure with periodicity enabled.}
The witness checks whether an early return would hide the error from the
renderer. Lean verifies that the over-budget index remains in the sampled
prefix under the modeled HDRx64 shadow arithmetic.

The fourth file uses \(c_\star=(-2+2^{-200})+0i\) at zoom \(2^{200}\)
and keeps the same output dimensions and sample density. The full
vertical span gives the standard request of 317 bits, which MPIR
rounds to 320 effective bits. At index 80, an ST replay with the
entered center loaded exactly has error about \(2.634\times10^{-62}\).
The quarter-sample tolerance is \(1.441\times10^{-64}\).
The calculator recommends 384 bits through that index.

For the GUI reproduction, keep the previous renderer and display
settings and select \emph{ST + Periodicity}. Clear the cached
references and enter the coordinates from the fourth file.
The file's \code{backend} selects the calculator's ST arithmetic.
The GUI's aspect fit increases the real center by \(2^{-321}\).
With respect to the exact orbit of that fitted center, the error at
index 80 is about \(1.140\times10^{-49}\). The Windows renderer
reports \code{ActualPeriodIfAny = 101}, so its reference includes
the sample at index 80.

The Lean search checks a finite family of real centers near \(-2\).
In its verified range, the real inputs and their square have MPIR
limb exponent one. The ST replay uses exact copies for the operations
with zero, trims the square's inputs and product at their respective
limb boundaries, and adds the negative parameter exactly.
The lower endpoint of the outward enclosure of the exact orbit
lies more than \(T\) above the replayed state at the witness index.
The theorem \code{PeriodicityWitness.checked\_failure} combines
the MPIR step model with bounds for the finite HDRx64 shadow operations.
For this prefix the reference shadow has magnitude between one and
two, and the derivative can grow by at most a factor \(11/2\) per step
relative to the larger of its previous magnitude and one.
With \(R_y=2^{-199}\), the near-return threshold stays below one
through index 80, and the reference bailout test also stays inactive.
The enabled schedule therefore retains a sample whose error exceeds
the budget. The kernel checks the numerical certificate for both
centers.

\subsection{MT9 Dynamic at the fixed viewport precision}
\label{subsec:reference-orbit-proof-mt9-viewport}

Dynamic MT9 can meet the viewport tolerance even if its internal state
differs from fixed MT9. With the same copied parameter
\(\widetilde c_\star\) and starting precision \(P\), the theorem
\code{Recommendation.reference\_error\_at} gives the fixed-orbit bound
\(F_n=2^{-P}J_n\) from \eqref{eq:reference-orbit-precision-estimator}.
The peak of the weights justifies the fixed-error induction, and the
weight at each index gives the budget used by the combined check.
By \code{MT9.disk\_prefix}, \(E_n\) bounds the dynamic-to-fixed
difference in the Euclidean norm, so each component difference is at most
\(E_n\). Consequently,
\begin{equation}
  \lVert z_n^{\mathrm D}-z_n^\star\rVert_1
  \le F_n+2E_n
  \le T,\qquad T=\eta\min(h_x,h_y).
  \label{eq:reference-orbit-proof-mt9-viewport}
\end{equation}
The theorem \code{MT9.viewport\_of\_recommendation} proves the sample
tolerance in \eqref{eq:reference-orbit-proof-sample-error} for the
published prefix when the combined budget holds at every positive
published index. The synthetic zero is exact. The numerical check uses
the calculator's weights and MT9's propagated enclosure, which can be
obtained without evaluating the fixed MPIR orbit. The theorem assumes the
MPIR and outward-arithmetic contracts and a fixed comparison prefix
covering the requested dynamic samples.

\code{MT9.fixed\_sufficiency\_alone} shows geometrically that a fixed
error at the boundary of the viewport budget can be pushed outside it by
an arbitrarily small positive drift. A bound above \(T\) leaves the
viewport tolerance undecided. The conversion tests in
\cref{subsec:reference-orbit-proof-mt9-certificate} concern renderer
inputs. The viewport budget can hold when conversion equivalence is
uncertain. Production MT9 logs conversion warnings and continues.

\begin{samepage}
\paragraph{Worked same-precision example.}
Use \((-2+2^{-200})+0i\) at zoom \(2^{200}\) with a \(5120\times2160\)
output, \(s=2\) (4x antialiasing), and \(\eta=1/4\). The calculator
returns 448 bits for fixed MT9 through index 99 and encodes the center
exactly. With the same center and MPIR precision, the native MT9 Dynamic
run with periodicity enabled lowers working precision to 384 bits.
The native check bounds the drift by \(2^{-216}\) at every positive
published index.
Lean checks that the peak fixed error plus twice that dyadic cap fits
within the quarter-sample budget. The result certifies the internal dynamic
prefix through index 99 at the starting precision selected for fixed MT9.

\end{samepage}

\subsubsection{MT9 Dynamic counterexample at 448 bits}
\label{subsec:reference-orbit-proof-mt9-counterexample}

Use the real center \(c_\star=-2+2^{-400}\), whose orbit stays near
the repelling point in \cref{subsec:reference-orbit-left-tip}.
Set zoom \(2^{450}\) on a \(1024\times1024\) output with \(s=1\)
and \(\eta=1/4\).
The quarter-sample tolerance is \(T=2^{-460}\).
Both MT9 runs start at a specified 448 bits with the same exactly
represented center. Both include iterations \(0\) through \(91\)
in the reference output with periodicity enabled. The first reduction
in the dynamic calculation occurs on the update from \(z_{90}\)
to \(z_{91}\), at 384 bits.

The proof in \code{MT9.RadiusWitness.counterexample} bounds the
fixed orbit's internal error by less than \(2^{-478}\) throughout
that prefix. At index 91, it bounds the dynamic orbit's error from
below by \(2^{-447}\), more than \(8192\) times the permitted error.
Both comparisons are against the exact mathematical orbit. The result
establishes that the actual accuracy of fixed MT9 at a given precision
can be sufficient for a radius that dynamic MT9 at the same starting
precision fails to meet. The theorem
\code{MT9.RadiusWitness.same\_radius\_failure} expresses the result
in sample spacings.

The Lean proof checks the real-axis update at the relevant limb
exponent. The MPIR copy can remove low limbs at a precision change,
and Karatsuba squares all remaining input limbs before trimming the
result. The aligned addition of the negative center then preserves
the exact sum on that grid. Lean replays the fixed and dynamic
arithmetic with exact dyadic operations and encloses the mathematical
orbit by rounding the interval endpoints outward. The native test
verifies the precision change, the number of samples in the
reference output, and agreement of the internal MPIR states with
the dyadic replay. The runtime conversion certificate accepts every
sample in the tested prefix. Those checks establish that the
references supplied for rendering agree at the stored precision.
The tighter viewport tolerance concerns bits of the internal MPIR
states that the stored samples do not retain.

For the same view, the calculator recommends 704 starting bits.
Both orbits meet the quarter-sample tolerance through index 91 at
that setting, as proved by
\code{MT9.RadiusWitness.\allowbreak meets\_at\_recommendation}.

\subsection{Periodicity and finite prefixes}
\label{subsec:reference-orbit-proof-periodicity}

Early termination controls how much reference data is generated and when
that reference can cover a later view. The schedule proofs identify
the last valid sample that precision and perturbation proofs can use,
including cases where the workers have already begun another update.

The exact derivative satisfies $d_{n+1}=2z_n^\star d_n+1$, with $d_0=0$
and $d_1=1$. The theorem \code{derivativeStep\_eq} proves that the real and
imaginary component updates satisfy this recurrence in exact arithmetic.
The Lean schedules model the comparison in
\cref{subsec:ref-orbit-periodicity} through the shadow inputs and the
comparison result for the observed index. In serial execution, a positive
near-return test stops after the current sample and before publishing a new
MPIR pair. In MT3, the caller evaluates the delayed shadow and honors the
stop as it collects the worker results, after publishing the delayed sample.
MT9 also publishes the delayed sample and completes its terminal
pipeline update before honoring a near return or bailout.
With periodicity disabled, the derivative is unchanged and the near-return
test cannot terminate the loop. Benchmark mode continues after a positive
test and retains the derivative for that body. The CPU exact and
finite-precision observation theorems hold in each of these modes.
The theorems concern the internal MPIR values associated with the observed
indices. The reduced-precision conversions and comparison outcomes are
inputs to the model, which applies the stopping rule for the selected mode.
